\section{House Robber}
\label{sec:dp-house-robber}

\problemstatement{Given an array $\textit{nums}$ where
$\textit{nums}[i]$ is the amount of money in house $i$, find the maximum
total amount that can be robbed without ever robbing two
\textbf{adjacent} houses.}

\begin{patternbox}
\emph{``Maximum sum, no two adjacent elements chosen''} is a different
recurrence \emph{shape} from every problem so far in this chapter: those
were all ``combine answers from a fixed set of earlier positions''
(summed, for counting, or minimized, for cost). This one is a genuine
\textbf{include-or-exclude choice} at every position -- structurally the
same choice as the pick/not-pick recursion from
the Recursion chapter's Power Set section, but now resolved by \textbf{maximizing}
rather than branching into both options forever.
\end{patternbox}

\subsection*{Deriving the recurrence}

Let $f(i)$ be the maximum amount robbable considering only houses $0$
through $i$. At house $i$, there are exactly two choices: \textbf{rob}
house $i$ (gaining $\textit{nums}[i]$, but then house $i-1$ \emph{must}
be skipped, so we add the best answer from $f(i-2)$), or \textbf{skip}
house $i$ entirely (carrying forward $f(i-1)$ unchanged, since skipping
loses nothing already secured). We take whichever choice is larger:
\[
  f(i) = \max\big(f(i-1),\ f(i-2) + \textit{nums}[i]\big).
\]
Base cases: $f(0) = \textit{nums}[0]$ (only one house, must take it for
the best possible answer at this point), and, treating an empty prefix
before index $0$ as contributing $0$, we can write $f(-1)=0$ for
notational convenience when computing $f(1) = \max(f(0), f(-1) +
\textit{nums}[1]) = \max(\textit{nums}[0], \textit{nums}[1])$.

\subsection*{Approach 1: Recursion}

\begin{lstlisting}
#include <bits/stdc++.h>
using namespace std;

int robRecursive(int i, vector<int>& nums) {
    if (i == 0) return nums[0];
    if (i < 0) return 0;

    int skip = robRecursive(i - 1, nums);
    int take = robRecursive(i - 2, nums) + nums[i];
    return max(skip, take);
}

int main() {
    vector<int> nums = {1, 2, 3, 1};
    int n = nums.size();
    cout << robRecursive(n - 1, nums) << endl; // 4
    return 0;
}
\end{lstlisting}

\begin{complexitybox}[Approach 1 Complexity]
\textbf{Time.} $O(2^n)$.
\textbf{Space.} $O(n)$ recursion stack.
\end{complexitybox}

\subsection*{Approach 2: Memoization}

\begin{lstlisting}
#include <bits/stdc++.h>
using namespace std;

int robMemo(int i, vector<int>& nums, vector<int>& memo) {
    if (i == 0) return nums[0];
    if (i < 0) return 0;
    if (memo[i] != -1) return memo[i];

    int skip = robMemo(i - 1, nums, memo);
    int take = robMemo(i - 2, nums, memo) + nums[i];
    return memo[i] = max(skip, take);
}

int main() {
    vector<int> nums = {1, 2, 3, 1};
    int n = nums.size();
    vector<int> memo(n, -1);
    cout << robMemo(n - 1, nums, memo) << endl; // 4
    return 0;
}
\end{lstlisting}

\begin{complexitybox}[Approach 2 Complexity]
\textbf{Time.} $O(n)$.
\textbf{Space.} $O(n)$ memo $+$ $O(n)$ recursion stack.
\end{complexitybox}

\subsection*{Approach 3: Tabulation}

\begin{lstlisting}
#include <bits/stdc++.h>
using namespace std;

int robTabulation(vector<int>& nums) {
    int n = nums.size();
    vector<int> dp(n);
    dp[0] = nums[0];

    for (int i = 1; i < n; i++) {
        int skip = dp[i - 1];
        int take = (i >= 2 ? dp[i - 2] : 0) + nums[i];
        dp[i] = max(skip, take);
    }
    return dp[n - 1];
}

int main() {
    vector<int> nums = {1, 2, 3, 1};
    cout << robTabulation(nums) << endl; // 4
    return 0;
}
\end{lstlisting}

\subsection*{Dry run of the tabulation table}

\dryrun{Take $\textit{nums} = [1,2,3,1]$.}

\begin{itemize}
  \item $\textit{dp}[0] = 1$.
  \item $\textit{dp}[1] = \max(\textit{dp}[0],\ 0+2) = \max(1,2) = 2$.
  \item $\textit{dp}[2] = \max(\textit{dp}[1],\ \textit{dp}[0]+3) =
        \max(2,\ 1+3) = \max(2,4) = 4$.
  \item $\textit{dp}[3] = \max(\textit{dp}[2],\ \textit{dp}[1]+1) =
        \max(4,\ 2+1) = \max(4,3) = 4$.
\end{itemize}

\begin{center}
\begin{tabular}{c|c|c|c|c}
$i$ & 0 & 1 & 2 & 3 \\ \hline
$\textit{nums}[i]$ & 1 & 2 & 3 & 1 \\ \hline
$\textit{dp}[i]$ & 1 & 2 & 4 & 4
\end{tabular}
\end{center}

The answer is $\textit{dp}[3]=4$, achieved by robbing houses $0$ and
$2$ (values $1+3=4$), which beats robbing house $1$ alone or house $3$
alone or houses $1,3$ together ($2+1=3$).

\begin{complexitybox}[Approach 3 Complexity]
\textbf{Time.} $O(n)$.
\textbf{Space.} $O(n)$ table, no recursion stack.
\end{complexitybox}

\subsection*{Approach 4: Space Optimization}

\begin{lstlisting}
#include <bits/stdc++.h>
using namespace std;

int robSpaceOptimized(vector<int>& nums) {
    int n = nums.size();
    int prev2 = 0, prev1 = nums[0];

    for (int i = 1; i < n; i++) {
        int skip = prev1;
        int take = prev2 + nums[i];
        int current = max(skip, take);
        prev2 = prev1;
        prev1 = current;
    }
    return prev1;
}

int main() {
    vector<int> nums = {1, 2, 3, 1};
    cout << robSpaceOptimized(nums) << endl; // 4
    return 0;
}
\end{lstlisting}

\begin{complexitybox}[Approach 4 Complexity]
\textbf{Time.} $O(n)$.
\textbf{Space.} $O(1)$.
\end{complexitybox}

\begin{takeawaybox}
``Maximum sum subject to a no-two-adjacent constraint'' is solved by an
include-or-exclude recurrence at every position, maximized rather than
summed -- the DP analogue of the pick/not-pick backtracking tree from
the Recursion chapter's Power Set section, but resolved in $O(n)$ instead of
$O(2^n)$ because we only need the \emph{best} outcome at each position,
not every possible outcome. Section~\ref{sec:dp-house-robber-2} adapts
this exact recurrence to a circular arrangement of houses.
\end{takeawaybox}
